\section{An injective nonsurjective cellular automaton}
\label{sec:cellular}

The scalar elements constructed in Section~\ref{sec:prescription} now
give a cellular automaton over the finite alphabet $K$. We prove the
standard transfer explicitly, in the right-neighbor convention used
in the introduction; compare\ \cite[Section~1]{ElekSzabo2004}.
The proof applies to all configurations, not just finitely supported
ones, and identifies the image by finite local equations.

\subsection{The group algebra acting on all configurations}

For this subsection, let $K$ be any field and $G$ any group, with
identity $1_G$. For $c=\sum_{u\in G}c_u u\in K[G]$, put
$\supp(c)=\{u\in G:c_u\ne0\}$. Multiplication in $K[G]$ is
\[
  (cd)_w=\sum_{uv=w}c_ud_v.
\]
Every sum has only finitely many nonzero terms. Define a $K$-linear map
$T_c:K^G\to K^G$ by
\begin{equation}
\label{eq:ca-operator}
  (T_cx)(g)=\sum_{u\in G}c_u x(gu).
\end{equation}
This sum is finite even when $x$ has infinite support. If $K$ is
finite, it is a cellular automaton with memory $\supp(c)$.

\begin{lemma}
\label{lem:ca-representation}
For every field $K$ and group $G$, the assignment $c\mapsto T_c$ is
a faithful unital $K$-algebra homomorphism from $K[G]$ to the algebra
of $K$-linear endomorphisms of $K^G$, with composition as multiplication.
In particular,
\begin{equation}
\label{eq:ca-composition}
  T_c\circ T_d=T_{cd},\qquad T_1=\id_{K^G}.
\end{equation}
\end{lemma}

\begin{proof}
Linearity in $c$ and the identity $T_1=\id_{K^G}$ follow directly
from Equation~\eqref{eq:ca-operator}. For $x\in K^G$ and $g\in G$,
\begin{align*}
 (T_c(T_dx))(g)
 &=\sum_{u\in G}c_u\sum_{v\in G}d_vx(guv)\\
 &=\sum_{w\in G}\left(\sum_{uv=w}c_ud_v\right)x(gw)
 =(T_{cd}x)(g).
\end{align*}
All sums being regrouped are finite. In particular, the product is
$cd$, not $dc$, for this right-neighbor convention.

For $h\in G$, let $\delta_h\in K^G$ take value $1_K$ at $h$ and
zero elsewhere. Then
\begin{equation}
\label{eq:ca-coefficient}
  (T_c\delta_h)(1_G)=c_h.
\end{equation}
Thus $T_c=0$ implies that every coefficient $c_h$ is zero, proving
faithfulness.
\end{proof}

\subsection{The local rule and its image}

\begin{proposition}
\label{prop:ca-transfer}
Let $K$ be a finite field, $G$ a group, and $a,b\in K[G]$ satisfy
$ab=1$. Put $M=\supp(b)$ and define
\begin{equation}
\label{eq:ca-rule}
  \phi:K^M\longrightarrow K,
  \qquad \phi(p)=\sum_{u\in M}b_up(u).
\end{equation}
The cellular automaton $\tau=T_b$ is injective, and
\begin{equation}
\label{eq:ca-image}
  \im(\tau)=\{y\in K^G:T_{ba}y=y\}.
\end{equation}
If $ba\ne1$, every $h\in G$ with $(ba-1)_h\ne0$ gives a
configuration $\delta_h\notin\im(\tau)$; hence $\tau$ is not
surjective.
\end{proposition}

\begin{proof}
The support of $b$ is finite, so Equation~\eqref{eq:ca-rule} is a
single finite-memory local rule, used at every $g\in G$. The alphabet
$K$ is finite and nonempty. Lemma~\ref{lem:ca-representation} gives
\begin{equation}
\label{eq:ca-leftinverse}
  T_a\circ\tau=T_{ab}=T_1=\id_{K^G}.
\end{equation}
Applying $T_a$ to $\tau(x)=\tau(x')$ yields $x=x'$ for arbitrary
$x,x'\in K^G$.

If $y=\tau(x)=T_bx$, then
\[
  T_{ba}y=T_bT_aT_bx=T_bx=y,
\]
by Equation~\eqref{eq:ca-leftinverse}. Conversely, if $T_{ba}y=y$,
then $y=T_b(T_ay)$ lies in the image of $\tau$. This proves
Equation~\eqref{eq:ca-image}.

If $(ba-1)_h\ne0$, Equation~\eqref{eq:ca-coefficient} gives
\[
  \bigl(T_{ba}\delta_h-\delta_h\bigr)(1_G)
  =\bigl(T_{ba-1}\delta_h\bigr)(1_G)
  =(ba-1)_h\ne0.
\]
Therefore $\delta_h$ is not fixed by $T_{ba}$, and
Equation~\eqref{eq:ca-image} excludes it from the image. This
calculation also covers $h=1_G$, when $\delta_h(1_G)=1_K$ rather
than zero.
\end{proof}

Equation~\eqref{eq:ca-image} describes the image by the uniform
finite local equations
\[
  \sum_{u\in\supp(ba-1)}(ba-1)_u\,y(gu)=0
  \qquad(g\in G).
\]
Thus the missing configuration has no preimage even among infinitely
supported functions. No restriction to finitely supported
configurations has entered the argument.

The memory need not be minimal. If a finite expression for $b$ is
$b=\sum_{j=1}^r\lambda_j u_j$, with $\lambda_j\in K$ and
$u_j\in G$ and with repetitions allowed, take
$M'=\{u_j:1\le j\le r\}$ and
$\phi'(p)=\sum_{j=1}^r\lambda_j p(u_j)$. Repeated group elements
automatically combine to give the same operator $T_b$. Consequently,
the finite word expressions of Equation~\eqref{eq:exact-outputs}
specify the local rule without an algorithm deciding which displayed
words represent the same group element.

\begin{corollary}
\label{cor:nonsurjunctive}
For the field $K$, group $G$, and outputs of
Theorem~\ref{thm:main}, the cellular automaton
$T_{b_{\mathrm{out}}}:K^G\to K^G$ is injective and not surjective.
Its image is the fixed set of $T_{b_{\mathrm{out}}a_{\mathrm{out}}}$.
In particular, $G$ is not surjunctive.
\end{corollary}

\begin{proof}
Apply Proposition~\ref{prop:ca-transfer} to
$a=a_{\mathrm{out}}$ and $b=b_{\mathrm{out}}$. Their product
relations are proved in Theorem~\ref{thm:main}, and the nonzero
reverse-product defect has a nonzero coefficient in the ordinary
group algebra.
\end{proof}

Characteristic two is used to construct the algebraic witnesses,
not in Proposition~\ref{prop:ca-transfer}. Finiteness of the field
supplies the finite alphabet, and the scalar form of the witnesses
makes that alphabet $K$ itself rather than a vector space over $K$.
The existence of a nonzero defect coefficient is certified
algebraically; no practical execution of the finite search or
normalization of group words is asserted.

\subsection{Transport to a group of type \texorpdfstring{$F_\infty$}{F infinity}}

A group has type $F_\infty$ if it admits a classifying CW complex with
finitely many cells in every dimension. The finite presentation of $G$
allows a further consequence: a separate embedding theorem transports
the same scalar counterexample to an overgroup of type $F_\infty$.

\begin{corollary}
\label{cor:finfty-counterexample}
For the finite characteristic-two field $K$ of Theorem~\ref{thm:main},
there are a group $U$ of type $F_\infty$ and scalar elements
$A,B\in K[U]$ such that $AB=1$ and $BA\ne1$.
The group $U$ is nonsofic and nonsurjunctive.
\end{corollary}

\begin{proof}
The group $G$ in Theorem~\ref{thm:main} is finitely presented.
The companion embedding theorem~\cite[Theorem~1.1]{UniversalFinftyCompanion}
therefore gives an injection $i:G\hookrightarrow U$ into a group of
type $F_\infty$. Extending $i$ linearly over the same field $K$ gives a
unital algebra injection $i_*:K[G]\hookrightarrow K[U]$: distinct group
basis elements remain distinct. Set $A=i_*(a_{\mathrm{out}})$ and
$B=i_*(b_{\mathrm{out}})$. Then
\[
 AB=1,\qquad
 BA-1=i_*(b_{\mathrm{out}}a_{\mathrm{out}}-1)\ne0.
\]
If $U$ were sofic, Elek--Szab\'o's theorem
\cite[Corollary~4.7]{ElekSzabo2004} would make $K[U]$ stably finite,
contradicting these scalar identities. Proposition~\ref{prop:ca-transfer},
applied to $U,K,A,B$, makes $T_B:K^U\to K^U$ an injective nonsurjective
cellular automaton over the finite alphabet $K$, with left inverse
$T_A$ in the same right-neighbor convention. Hence $U$ is nonsurjunctive.
\end{proof}
